% basic nastaveni
\documentclass[12pt]{article}
\setlength{\parindent}{0pt}
\setlength{\parskip}{1em}
\usepackage[utf8]{inputenc}
\usepackage[czech]{babel}
\usepackage[T1]{fontenc}
\usepackage{graphics}
\usepackage{caption}
\usepackage{hyperref}
\usepackage{xcolor}
\usepackage{float}

% matematicke package
\usepackage{amsmath}
\usepackage{amsfonts}
\usepackage{amssymb}
\usepackage{geometry}
\geometry{margin=2.5cm}
\usepackage{pgfplots}
\pgfplotsset{compat=1.18}
\usepackage{mathrsfs}

\title{Maturitní otázka číslo 15: Tělesa}
\author{}
\date{}

\begin{document}
\maketitle
\subsection*{Co to jsou tělesa}
Prostorově omezené souvislé útvary se nazývají tělesa. Jejich hranicí jsou uzavřené plochy. Vnitřní
body jsou body, které se nacházejí uvnitř prostoru. Hraniční body se nacházejí na plochách daného
tělesa. Každé těleso musí mít nenulový objem. Pokud všechny úsečky tvořené všemi dvojicemi bodů
tělesa leží uvnitř tělesa, jde o konvexní těleso, jinak jde o nekonvexní.
\subsection*{Mnohostěny}
Mnohostěny se skládají z vrcholů, hran a stěn. Vrcholy jsou body, kde se stýkají hrany, hrany jsou
úsečky, kde se stýkají dvě stěny a stěny jsou mnohoúhelníky tvořící povrch tělesa. Dále v nich
rozlišujeme stěnové a tělesové úhlopříčky. Stěnové úhlopříčky spojují vrcholy vrámci jedné stěny,
zatímco tělesové spojují vrcholy, které neleží v téže stěně.
\subsubsection*{Hranoly}
Hranoly jsou tělesa, které mají dvě shodné podstavy v rovnoběžných rovinách a boční stěny, které
jsou tvořeny rovnoběžníky. Hranoly rozlišujeme na n-boké podle toho, kolik hran má jejích podstava.
Rozlišujeme dva specifické hranoly a to kolmý hranol, jehož boční stěny jsou kolmé k podstavám, takže
všechny boční stěny jsou obdelníky nebo čtverce. Dalším z nich je pravidelný hranol, což je typ
kolmého hranolu, jehož podstavou je pravidelný n-úhelník. U pravidelného hranolu platí, že délka
všech tělesových úhlopříček je stejná.

$V = S_p \cdot v$, $S_p$ je obsah podstavy a $v$ je výška hranolu

$S = 2 \cdot S_p + S_{pl}$, $S_{pl}$ je obsah pláště.

Kvádr je speciální typem hranolu. Pokud je jeho podstavou obdélník, jde o kolmý čtýřboký hranol, který
má délku podstav $a, b$ a výšku $c$. Speciálním typem kvádru je čtvercový kvádr, který má jako 
podstavu čtverec je to pravidelný čtyřboký hranol. Kvádr má čtyři tělesové úhlopříčky, které jsou
stejně dlouhé a jejich délka lze vypočítat jako $u_t = \sqrt{a^2 + b^2 + c^2}$.

$V = a \cdot b \cdot c$

$S = 2(a \cdot b + b \cdot c + a \cdot c)$

Krychle je pravidelný čtyřboký hranol, který má jako podstavu čtverec délky $a$ a výšku taky $a$.
Všechny tělesové úhlopříčky mají stejnou délku $u_t = a \cdot \sqrt{3}$ a všechny stěnové $u_s =
a \cdot \sqrt{2}$.

$V = a^3$

$S = 6 \cdot a^2$

\subsection*{Jehlany}
\end{document}
